% !TEX root = ../main.tex
\chapter{Conclusiones}

Con respecto al primer objetivo específico, se seleccionó UBC-OCEAN y se
derivaron una población no supervisada de parches, la clasificación de cinco
subtipos a nivel de WSI y la segmentación de tumor, estroma y necrosis. Las
particiones por WSI mantuvieron separados el entrenamiento, la validación y la
prueba. Las máscaras permitieron evaluar la segmentación a resolución de
parche, pero no de forma exhaustiva porque solo anotan parte de cada lámina.

En cumplimiento del segundo objetivo, se implementaron dos VAE residuales
comparables: uno ordinario y otro equivariante a \(\mathrm{SO}(2)\) sobre
campos \(F_0\) y \(F_1\). Ambos conservaron los bloques, no linealidades y
variables latentes, y alimentaron los mismos clasificadores. La rama
equivariante empleó menos parámetros libres, pero soportes mayores y más
cómputo estimado; se compararon, por tanto, dos arquitecturas completas y no
un experimento en el que solo cambiara una variable.

En relación con el tercer objetivo, la hipótesis de que la representación
equivariante produciría una mejora general en reconstrucción, clasificación y
segmentación no pudo validarse con evidencia estadística concluyente. El VAE
base obtuvo estimaciones puntuales ligeramente mejores en las métricas de
reconstrucción. El VAE-\(\mathrm{SO}(2)\) obtuvo mejores valores puntuales de
macro-F1 y exactitud en
diagnóstico, y mayor macro-F1 en cuatro de los cinco presupuestos reducidos de
etiquetas de tejido; sin embargo, los intervalos globales de estas
comparaciones incluyeron cero y solo el presupuesto de 500 etiquetas por clase
tuvo un intervalo simultáneo que lo excluyó. La inversión de la estimación
puntual al usar las \(5\,671\) etiquetas disponibles sugiere que cualquier
beneficio puede concentrarse en el régimen de supervisión muy escasa, pero los
datos actuales no permiten afirmarlo como una ventaja uniforme de eficiencia
muestral.

Por otra parte, el resultado más claro apareció en la geometría del espacio
latente. Para cuartos de vuelta exactos, la mediana de la razón de error de
acción disminuyó de 0,709453 en el VAE base a \(2,18\times10^{-6}\) en el
VAE-\(\mathrm{SO}(2)\), y la transformación del latente produjo reconstrucciones
visualmente coherentes sin los artefactos observados en la rama ordinaria. Las
razones de acción y canonización también fueron menores bajo rotaciones
interpoladas, aunque no alcanzaron todos los umbrales predefinidos. La
arquitectura aprendió y conservó, por tanto, una regularidad rotacional
controlable en la representación; este resultado no implica todavía una
factorización global del espacio latente ni una equivarianza continua perfecta
de extremo a extremo.

En síntesis, el presente trabajo produjo un modelo semisupervisado funcional y
evaluó sus propiedades predictivas y geométricas. Su principal contribución no
consiste en demostrar superioridad general, sino en mostrar que una restricción
continua de rotación puede conservarse de manera mensurable en la
representación latente. Una posible ventaja bajo pocas etiquetas permanece
como hipótesis y requiere más evidencia.
